\begin{abstract}
Researchers in electrical and computer engineering increasingly delegate literature
search, coding, data analysis and drafting to AI agents, yet practices for supervising
several agents at once lag behind adoption, and the researcher remains accountable for
every claim. Existing multi-agent frameworks and levels-of-automation models do not tell a
researcher how much autonomy to grant a task, how work should pass between agents, or how
many agents one person can supervise. This paper presents an orchestration framework that
answers these questions. Each task's autonomy level is set from three properties of its
effects (reversibility, reach and verifiability) by a rule that is complete and
monotone. Work passes between agents as version-controlled files that no agent may mark
as accepted when it produced them, and every reported number must name the computation
that produced it. A review-load model bounds team size by the researcher's daily
attention and shows that the ratio of planning to execution agents follows from their
review costs rather than being fixed. We apply the framework to an illustrative
analog-to-digital converter characterization task and to two AI-assisted publication
projects, a 20-chapter textbook and an arXiv paper, whose logs and repositories record 13
defects. Mechanical checks caught every defect for which one existed, errors of meaning
were caught only when someone recomputed or read the work, and a quality score that no
computation had produced escaped every control and was published. The results suggest
measuring review load and escaped defects before enlarging an agent team.
\end{abstract}

\noindent\textbf{Keywords:} AI agents, human supervisory control, levels of autonomy,
multi-agent systems, research workflows, verification, reproducibility, electrical
engineering
